\subsection{EXPANSION OF THE REAL EXPONENTIAL}

Since \(\mathrm{D}^n (e^x) = e^x\) the Taylor expansion of order \(n\) for \(e^x\) is

\[
e ^ {x} = 1 + \frac {x}{1 !} + \frac {x ^ {2}}{2 !} + \dots + \frac {x ^ {n}}{n !} + \int_ {0} ^ {x} \frac {(x - t) ^ {n}}{n !} e ^ {t} d t.\tag{1}
\]

The remainder in this formula is \(>0\) for \(x > 0\) and has the sign of \((-1)^{n + 1}\) when \(x < 0\) ; moreover, the inequality of the mean shows that

\[
\frac {x ^ {n + 1}}{(n + 1) !} <   \int_ {0} ^ {x} \frac {(x - t) ^ {n}}{n !} e ^ {t} d t <   \frac {x ^ {n + 1} e ^ {x}}{(n + 1) !} \quad \text { for } x > 0\tag{2}
\]

\[
\frac {\left| x ^ {n + 1} \right| e ^ {x}}{(n + 1) !} <   \left| \int_ {0} ^ {x} \frac {(x - t) ^ {n}}{n !} e ^ {t} d t \right| <   \frac {\left| x ^ {n + 1} \right|}{(n + 1) !} \quad \text { for } x <   0\tag{3}
\]

Now one knows that the sequence \((x^{n}/n!)\) has limit 0 when n increases indefinitely, for all \(x \geqslant 0\) (Gen.~Top., IV, p.~365); thus, keeping x fixed and letting n grow indefinitely in (1) it follows, from (2) and (3), that

\[
e ^ {x} = \sum_ {n = 0} ^ {\infty} \frac {x ^ {n}}{n !}\tag{4}
\]

and the series on the right-hand side is absolutely and uniformly convergent on every compact interval in \(\mathbf{R}\) . In particular, one has the formula

\[
e = 1 + \frac {1}{1 !} + \frac {1}{2 !} + \dots + \frac {1}{n !} + \dots\tag{5}
\]

This formula allows us to calculate rational approximations as close as we desire to the number \(e\) ; one obtains

\[
e = 2. 7 1 8   2 8 1   8 2 8 \dots
\]

to within \(1/10^{9}\) . Formula (5) proves, moreover, that e is an irrational number \(^{2}\) (Gen.~Top., IV, p.~375).

Remark. Since the remainder in formula (1) is \textgreater0 for x \textgreater{} 0 one has, for x \textgreater{} 0

\[
e ^ {x} > 1 + \frac {x}{1 !} + \frac {x ^ {2}}{2 !} + \dots + \frac {x ^ {n + 1}}{(n + 1) !}
\]

and a fortiori

\[
e ^ {x} > \frac {x ^ {n + 1}}{(n + 1) !}
\]

for every integer \(n\) : one deduces from this that \(e^x / x^n\) tends to \(+\infty\) with \(x\) , for every integer \(n\) : we shall find this result again in chap.~V by another method (V, p.~231).

